% answers_short.tex —— 2007 年·普通高等学校招生考试·四川卷·文科 · 简短答案速查
% 由 tools/make_short_answers.py 自 answers.tex 的 \gksolhead/\gkanswer 提取，
% 并与 2007/2007_数学真题_四川文_答案要点.md 逐题比对后生成，勿手改（改源 MD 后重跑）。

\begin{gknotice}
\textbf{使用说明}：本节只给最终答案便于快速核对，完整推导见解析版。答案为编写者独立推导，\textbf{非官方标准答案}。
\end{gknotice}

\begin{gknotice}
选择题答案中的 U 表示待核，未给出猜测答案，详情见解析版。
\end{gknotice}

\gkpart{选择题}
\begin{center}\begin{tabular}{ccccccccccccc}
\toprule
题号 & 1 & 2 & 3 & 4 & 5 & 6 & 7 & 8 & 9 & 10 & 11 & 12 \\
\midrule
答案 & A & C & B & D & A & C & B & A & B & C & B & U \\
\bottomrule
\end{tabular}\end{center}

\gkpart{填空题}
\begin{itemize}[leftmargin=2.6em]
  \item[\textbf{13.}] $8$。
  \item[\textbf{14.}] $\frac{\pi}{6}$。
  \item[\textbf{15.}] $x=\frac32$。
  \item[\textbf{16.}] ①④。
\end{itemize}

\gkpart{解答题}
\begin{itemize}[leftmargin=2.6em]
  \item[\textbf{17.}] （1）$0.7599$；（2）不合格 $1$ 件的概率为 $\frac{51}{190}$，不合格 $2$ 件的概率为 $\frac{3}{190}$，拒收概率为 $\frac{27}{95}$。
  \item[\textbf{18.}] （1）$\tan2\alpha=-\frac{8\sqrt3}{47}$；（2）$\beta=\frac{\pi}{3}$。
  \item[\textbf{19.}] （1）证明见解析；（2）$\arccos\frac{\sqrt{39}}{13}$（约 $61.3^\circ$）；（3）$\frac{\sqrt6}{6}$。
  \item[\textbf{20.}] （1）$a=2$，$b=-12$，$c=0$；（2）递增区间为 $(-\infty,-\sqrt2)$ 与 $(\sqrt2,+\infty)$，在 $[-1,3]$ 上最大值为 $18$、最小值为 $-8\sqrt2$。
  \item[\textbf{21.}] （1）$\left(1,\frac{\sqrt3}{2}\right)$；（2）$\left(-2,-\frac{\sqrt3}{2}\right)\cup\left(\frac{\sqrt3}{2},2\right)$。
  \item[\textbf{22.}] （1）$x_{n+1}=\frac{x_n^2+4}{2x_n}$；（2）$\{a_n\}$ 是首项为 $\lg3$、公比为 $2$ 的等比数列，$x_n=\frac{2\left(3^{2^{n-1}}+1\right)}{3^{2^{n-1}}-1}$；（3）证明见解析。
\end{itemize}
